\documentclass[11pt]{article}
\usepackage[margin=1in]{geometry}
\usepackage{amsmath}

\setlength{\parindent}{0pt}

\begin{document}

\textbf{Question 1:}
\bigskip

Indicate whether each of the following actions will increase or decrease a bond's yield to maturity:

\begin{enumerate}
  \item[a.] The bond's price increases. 
  \item[b.] The bond is downgraded by the rating agencies. 
  \item[c.] A change in the bankruptcy code makes it more difficult for bondholders to receive payments in the event the firm declares bankruptcy.
  \item[d.] The economy seems to be shifting from a boom to a recession. Discuss the effects of the firm's credit strength in your answer.
  \item[e.] Investors learn that the bonds are subordinated to another debt issue.
\end{enumerate}
 \bigskip

\textbf{Answer}

\begin{enumerate}
  \item[a.] Decrease 
  \item[b.] Increase
  \item[c.] Increase 
  \item[d.] Increase. A recession would increase the default risk of the bond and decrease the credit strength of the security. 
  \item[e.] Increase
\end{enumerate}
\bigskip
\bigskip

\textbf{Question 2:}
\bigskip

A bond has a \$1,000 par value, 12 years to maturity, and an 8\% annual coupon and sells for \$980. 

\begin{enumerate}
  \item[a.] What is its yield to maturity (YTM)? 
  \item[b.] Assume that the yield to maturity remains constant for the next three years. What will the price be 3 years from today?
\end{enumerate}
\bigskip

\textbf{Answer}

\begin{enumerate}
  \item[a.]
    \[
      980 = \sum_{i=1}^{12} \frac{80}{(1+YTM)^i} + \frac{1000}{(1+YTM)^{12}}
    \]
    \begin{center}
    After plugging into excel: YTM \approx \boxed{8.27\%}
    \end{center}

  \item[b.]
    \[
      P = \sum_{i=1}^{9} \frac{80}{(1+YTM)^i} + \frac{1000}{(1+YTM)^{9}}
    \]
    \[
      \Rightarrow P = 494.19 + 489.13 = \boxed{\$983.32}
    \]
\end{enumerate}

\end{document}